\chapter{Généralités sur l'industrie automobile}
\label{cp:generalites-automobile}

\textit{Ce deuxième chapitre vise à établir un socle de connaissances générales sur le secteur automobile, en abordant ses fondements historiques, ses évolutions technologiques et les grands enjeux auxquels il fait face aujourd'hui. Il permettra de situer le contexte industriel dans lequel s'inscrit ce projet, en mettant en lumière les principales étapes du cycle de développement d'un véhicule, les exigences réglementaires en matière de sécurité, ainsi que les approches numériques de plus en plus sollicitées pour anticiper et optimiser les performances des véhicules. Ces notions fondamentales serviront de base pour la compréhension des travaux réalisés dans la suite de ce rapport.}

\section{Classification des segments automobiles}

Dans l'industrie automobile, les véhicules sont classés par segments selon leurs dimensions, leur type d'utilisation et leur positionnement sur le marché. Cette classification permet aux constructeurs de cibler différents profils de clientèle et de structurer leur gamme de produits. Le \autoref{tab:segments-automobiles} présente les principaux segments utilisés en Europe.

\begin{table}[H]
    \caption{Classification des segments automobiles}
    \label{tab:segments-automobiles}
    \centering
    \small
    \renewcommand{\arraystretch}{1.4}
    \begin{tabularx}{\textwidth}{|>{\cellcolor{ensablue!15}\bfseries\centering\arraybackslash}p{1.8cm}|>{\centering\arraybackslash}X|>{\centering\arraybackslash}p{2.8cm}|>{\centering\arraybackslash}X|}
        \hline
        \rowcolor{ensablue!25}
        \textbf{Segment} & \textbf{Dénomination} & \textbf{Dimensions} & \textbf{Type d'utilisation} \\
        \hline
        B0 & Micro Urbaine & < 3,1 m & Zone urbaine -- courte distance \\
        \hline
        A & Urbaine & 3,1 m -- 3,6 m & Zone urbaine -- courte et moyenne distance \\
        \hline
        B & Citadine polyvalente & 3,6 m -- 4,1 m & Multi utilisations \\
        \hline
        B+ & Citadine monospace & 3,6 m -- 4,1 m & Multi utilisations \\
        \hline
        C & Compactes & 4,1 m -- 4,5 m & Familiale -- longues distances \\
        \hline
        D & Berlines & 4,5 m -- 4,8 m & Familiale confortable -- longue distance \\
        \hline
        E & Grandes routières & 4,9 m -- 5 m & Familiale confortable -- longue distance \\
        \hline
        F & Berlines de luxe & 5 m ou plus & Voiture luxueuse \\
        \hline
        SUV -- Crossover & Véhicule utilitaire sportif & 4,1 m -- 5 m & Familiale confortable -- longue distance -- parfois 4$\times$4 \\
        \hline
    \end{tabularx}
\end{table}
